\chapter{L14 -- Forma standard di raggiungibilità, forma canonica di controllo e introduzione all'osservabilità}

\section{Obiettivo della lezione}

In questa lezione chiudiamo il blocco sulla raggiungibilità e iniziamo quello sull’osservabilità.

Più precisamente:
\begin{itemize}
    \item usiamo un esempio per costruire esplicitamente la \textbf{forma standard di raggiungibilità};
    \item vediamo quando conviene passare alla \textbf{forma canonica di controllo};
    \item capiamo che gli \textbf{autovalori non raggiungibili} non possono essere spostati con una retroazione di stato;
    \item introduciamo il problema duale dell’\textbf{osservabilità}, cioè la ricostruzione dello stato a partire da ingressi e uscite misurate.
\end{itemize}

\bigskip

Il filo logico è questo:
\[
\text{raggiungibilità} \quad \Longleftrightarrow \quad \text{quali componenti dello stato posso muovere con } u,
\]
\[
\text{osservabilità} \quad \Longleftrightarrow \quad \text{quali componenti dello stato posso ricostruire da } (u,y).
\]

%=========================================================
\section{Esempio guida: forma standard di raggiungibilità}

Consideriamo il sistema discreto
\[
x(k+1)=Ax(k)+Bu(k),\qquad y(k)=Cx(k),
\]
con
\[
A=
\begin{bmatrix}
0 & 0 & 2\\
0 & \dfrac{\alpha}{10} & 0\\
0 & 0 & 0
\end{bmatrix},
\qquad
B=
\begin{bmatrix}
0\\
0\\
1
\end{bmatrix},
\qquad
C=
\begin{bmatrix}
0 & 1 & 1
\end{bmatrix},
\]
dove \(\alpha\in\{0,1,\dots,9\}\).

\subsection{Autovalori e struttura di Jordan}

Gli autovalori di \(A\) sono
\[
\lambda_1=0,\qquad \lambda_2=0,\qquad \lambda_3=\frac{\alpha}{10}.
\]

Quindi ci sono due casi distinti.

\subsubsection*{Caso 1: \(\alpha\neq 0\)}

In questo caso:
\[
m_a(0)=2,\qquad m_a\!\left(\frac{\alpha}{10}\right)=1.
\]

Poiché il blocco in alto a destra è non nullo, i due autovalori nulli non corrispondono a due autovettori indipendenti: infatti
\[
m_g(0)=\dim\ker(A)=1.
\]
Dunque l’autovalore \(0\) è associato a \textbf{un unico blocco di Jordan di ordine 2}.

Per l’autovalore \(\alpha/10\neq 0\) si ha invece
\[
m_g\!\left(\frac{\alpha}{10}\right)=1,
\]
quindi esso corrisponde a un blocco \(1\times 1\).

In sintesi, per \(\alpha\neq 0\), la forma di Jordan è simile a
\[
J=
\begin{bmatrix}
0 & 1 & 0\\
0 & 0 & 0\\
0 & 0 & \dfrac{\alpha}{10}
\end{bmatrix}.
\]

\subsubsection*{Caso 2: \(\alpha=0\)}

Se \(\alpha=0\), allora l’unico autovalore è
\[
\lambda=0,\qquad m_a(0)=3.
\]
Inoltre
\[
m_g(0)=\dim\ker(A)=2,
\]
perciò ci sono \textbf{due blocchi di Jordan} associati all’autovalore \(0\): uno di ordine \(2\) e uno di ordine \(1\).

Una possibile forma di Jordan è
\[
J=
\begin{bmatrix}
0 & 1 & 0\\
0 & 0 & 0\\
0 & 0 & 0
\end{bmatrix}.
\]

\subsection{Matrice di raggiungibilità}

Calcoliamo
\[
R_3=\begin{bmatrix}B & AB & A^2B\end{bmatrix}.
\]

Poiché
\[
B=
\begin{bmatrix}
0\\0\\1
\end{bmatrix},
\qquad
AB=
\begin{bmatrix}
2\\0\\0
\end{bmatrix},
\qquad
A^2B=A(AB)=
\begin{bmatrix}
0\\0\\0
\end{bmatrix},
\]
si ottiene
\[
R_3=
\begin{bmatrix}
0 & 2 & 0\\
0 & 0 & 0\\
1 & 0 & 0
\end{bmatrix}.
\]

Dunque
\[
\operatorname{Im}(R_3)=\operatorname{span}\!\left\{
\begin{bmatrix}0\\0\\1\end{bmatrix},
\begin{bmatrix}1\\0\\0\end{bmatrix}
\right\},
\qquad
\operatorname{rank}(R_3)=2.
\]

Ne segue che il sistema \textbf{non è completamente raggiungibile}, perché lo spazio di stato ha dimensione \(3\) ma lo spazio raggiungibile ha dimensione \(2\).

\subsection{Catena degli spazi raggiungibili}

Si ha
\[
\mathcal R_1=\operatorname{Im}(B)=
\operatorname{span}\!\left\{
\begin{bmatrix}0\\0\\1\end{bmatrix}
\right\},
\]
e
\[
\mathcal R_2=\operatorname{Im}\!\begin{bmatrix}B & AB\end{bmatrix}
=
\operatorname{span}\!\left\{
\begin{bmatrix}0\\0\\1\end{bmatrix},
\begin{bmatrix}1\\0\\0\end{bmatrix}
\right\}.
\]
Poiché \(A^2B=0\), si ha già
\[
\mathcal R_2=\mathcal R_3.
\]

Quindi la parte raggiungibile è il piano
\[
\mathcal R=\operatorname{span}\{e_1,e_3\}.
\]

\subsection{Scelta di una base adattata}

Scegliamo come base di \(\mathcal R\)
\[
T_R=\begin{bmatrix}e_1 & e_3\end{bmatrix},
\]
e completiamola con un vettore della parte non raggiungibile, ad esempio
\[
T_N=\begin{bmatrix}e_2\end{bmatrix}.
\]
Allora una matrice di cambio di base è
\[
T=\begin{bmatrix}T_R & T_N\end{bmatrix}
=
\begin{bmatrix}
1 & 0 & 0\\
0 & 0 & 1\\
0 & 1 & 0
\end{bmatrix}.
\]
Essendo una matrice di permutazione, vale
\[
T^{-1}=T^\top.
\]

Passando alle nuove coordinate \(x=Tz\), otteniamo
\[
\widetilde A=T^{-1}AT,
\qquad
\widetilde B=T^{-1}B,
\qquad
\widetilde C=CT.
\]

Un calcolo diretto fornisce
\[
\widetilde A=
\begin{bmatrix}
0 & 2 & 0\\
0 & 0 & 0\\
0 & 0 & \dfrac{\alpha}{10}
\end{bmatrix},
\qquad
\widetilde B=
\begin{bmatrix}
0\\1\\0
\end{bmatrix},
\qquad
\widetilde C=
\begin{bmatrix}
0 & 1 & 1
\end{bmatrix}.
\]

Questa è esattamente la \textbf{forma standard di raggiungibilità}:
\[
\widetilde A=
\begin{bmatrix}
A_R & A_{RN}\\
0 & A_N
\end{bmatrix},
\qquad
\widetilde B=
\begin{bmatrix}
B_R\\
0
\end{bmatrix}.
\]

Nel nostro caso:
\[
A_R=
\begin{bmatrix}
0 & 2\\
0 & 0
\end{bmatrix},
\qquad
A_N=
\begin{bmatrix}
\dfrac{\alpha}{10}
\end{bmatrix}.
\]

Quindi:
\begin{itemize}
    \item la dinamica su \(z_R\) è la parte raggiungibile;
    \item la dinamica su \(z_N\) è la parte non raggiungibile.
\end{itemize}

\subsection{Interpretazione degli autovalori}

In questa decomposizione:
\begin{itemize}
    \item gli autovalori di \(A_R\) sono gli \textbf{autovalori raggiungibili};
    \item gli autovalori di \(A_N\) sono gli \textbf{autovalori non raggiungibili}.
\end{itemize}

Nel nostro esempio:
\begin{itemize}
    \item l’autovalore \(0\) associato al blocco \(A_R\) è raggiungibile;
    \item l’autovalore \(\alpha/10\) è non raggiungibile.
\end{itemize}

Se \(\alpha=0\), allora anche uno dei modi nulli è non raggiungibile: infatti un autovalore può comparire sia nella parte raggiungibile sia nella parte non raggiungibile.

%=========================================================
\section{Verifica rapida con PBH}

Il lemma di Popov--Belevitch--Hautus (PBH) dice che la coppia \((A,B)\) è completamente raggiungibile se e solo se
\[
\operatorname{rank}\!\begin{bmatrix}A-\lambda I & B\end{bmatrix}=n
\qquad
\forall \lambda\in\mathbb C.
\]

Nel nostro esempio, l’autovalore sospetto è
\[
\lambda=\frac{\alpha}{10}.
\]
Infatti
\[
A-\frac{\alpha}{10}I=
\begin{bmatrix}
-\frac{\alpha}{10} & 0 & 2\\
0 & 0 & 0\\
0 & 0 & -\frac{\alpha}{10}
\end{bmatrix},
\]
e quindi
\[
\begin{bmatrix}
A-\frac{\alpha}{10}I & B
\end{bmatrix}
=
\begin{bmatrix}
-\frac{\alpha}{10} & 0 & 2 & 0\\
0 & 0 & 0 & 0\\
0 & 0 & -\frac{\alpha}{10} & 1
\end{bmatrix}.
\]
Il rango è \(2\), non \(3\), dunque \(\alpha/10\) è effettivamente \textbf{non raggiungibile}.

\bigskip

Questa verifica conferma quanto avevamo già visto dalla forma standard di raggiungibilità.

%=========================================================
\section{Ricostruzione di \(x(0)\) dalle uscite: un esempio in tempo discreto}

Restiamo nello stesso esempio, e fissiamo
\[
\alpha=5,
\qquad
u(0)=u(1)=-1.
\]
Supponiamo di volere trovare gli stati iniziali \(x(0)\) tali che
\[
y(0)=3,\qquad y(1)=1,\qquad y(2)=0.
\]

\subsection{Espansione delle uscite}

Poiché il sistema è discreto e \(D=0\), si ha
\[
y(k)=Cx(k).
\]
Inoltre
\[
x(1)=Ax(0)+Bu(0),
\]
\[
x(2)=A^2x(0)+ABu(0)+Bu(1).
\]
Di conseguenza
\[
y(0)=Cx(0),
\]
\[
y(1)=CAx(0)+CB\,u(0),
\]
\[
y(2)=CA^2x(0)+CAB\,u(0)+CB\,u(1).
\]

Portando a sinistra i termini noti legati agli ingressi, otteniamo
\[
\begin{bmatrix}
y(0)\\
y(1)-CBu(0)\\
y(2)-CABu(0)-CBu(1)
\end{bmatrix}
=
\begin{bmatrix}
C\\
CA\\
CA^2
\end{bmatrix}
x(0).
\]

\subsection{Calcolo esplicito}

Per \(\alpha=5\), cioè \(\alpha/10=1/2\), si ha
\[
A=
\begin{bmatrix}
0 & 0 & 2\\
0 & \frac12 & 0\\
0 & 0 & 0
\end{bmatrix},
\qquad
C=
\begin{bmatrix}
0 & 1 & 1
\end{bmatrix},
\qquad
B=
\begin{bmatrix}
0\\0\\1
\end{bmatrix}.
\]
Calcoliamo:
\[
CB=
\begin{bmatrix}
0 & 1 & 1
\end{bmatrix}
\begin{bmatrix}
0\\0\\1
\end{bmatrix}=1,
\]
\[
CAB=CAB=0.
\]
Quindi il vettore corretto delle uscite è
\[
\Delta y=
\begin{bmatrix}
3\\
1-1\cdot(-1)\\
0-0\cdot(-1)-1\cdot(-1)
\end{bmatrix}
=
\begin{bmatrix}
3\\2\\1
\end{bmatrix}.
\]

La matrice di osservabilità ai primi tre istanti è
\[
\mathcal O=
\begin{bmatrix}
C\\CA\\CA^2
\end{bmatrix}
=
\begin{bmatrix}
0 & 1 & 1\\
0 & \frac12 & 0\\
0 & \frac14 & 0
\end{bmatrix}.
\]
Dobbiamo quindi risolvere
\[
\begin{bmatrix}
0 & 1 & 1\\
0 & \frac12 & 0\\
0 & \frac14 & 0
\end{bmatrix}
x(0)=
\begin{bmatrix}
3\\2\\1
\end{bmatrix}.
\]

Se scriviamo
\[
x(0)=
\begin{bmatrix}
x_1(0)\\x_2(0)\\x_3(0)
\end{bmatrix},
\]
le equazioni diventano
\[
x_2(0)+x_3(0)=3,
\]
\[
\frac12 x_2(0)=2,
\]
\[
\frac14 x_2(0)=1.
\]
Le ultime due sono coerenti e forniscono
\[
x_2(0)=4.
\]
Sostituendo nella prima:
\[
x_3(0)=3-4=-1.
\]
La variabile \(x_1(0)\) non compare mai, quindi resta libera.

Conclusione:
\[
x(0)=
\begin{bmatrix}
\xi\\
4\\
-1
\end{bmatrix},
\qquad \xi\in\mathbb R.
\]

\subsection{Interpretazione}

Questo esempio è molto istruttivo:
\begin{itemize}
    \item le misure di uscita permettono di ricostruire \(x_2(0)\) e \(x_3(0)\);
    \item non permettono invece di ricostruire \(x_1(0)\).
\end{itemize}

Quindi il sistema \emph{non è completamente osservabile}: esiste almeno una direzione dello stato che non lascia traccia sull’uscita.

%=========================================================
\section{Osservazione strutturale: autovalori multipli e numero minimo di ingressi}

Un punto importante, che negli appunti compare come osservazione sui blocchi di Jordan, è il seguente.

\subsection*{Osservazione}
Se un autovalore \(\lambda\) è associato a \(\nu\) blocchi di Jordan distinti, allora la sua molteplicità geometrica è
\[
m_g(\lambda)=\nu.
\]

Perché la coppia \((A,B)\) possa essere completamente raggiungibile, la matrice
\[
\begin{bmatrix}
A-\lambda I & B
\end{bmatrix}
\]
deve avere rango massimo.

Ma in forma di Jordan, il blocco \(J-\lambda I\) perde rango esattamente di \(\nu\), cioè di tanti quanti sono i blocchi di Jordan associati a \(\lambda\). Per recuperare quel rango servono abbastanza colonne in \(B\) che agiscano in modo indipendente sulle direzioni giuste.

Ne segue la regola pratica:

\begin{quote}
\textbf{se per un autovalore \(\lambda\) ci sono \(\nu\) blocchi di Jordan, allora servono almeno \(\nu\) ingressi indipendenti per avere la possibilità di rendere raggiungibile quel modo.}
\end{quote}

In particolare, se il numero di ingressi è troppo piccolo rispetto al numero di blocchi associati a uno stesso autovalore, la completa raggiungibilità è impossibile.

Questa osservazione non sostituisce PBH, ma lo rende intuitivo.

%=========================================================
\section{Se \((A,B)\) non è completamente raggiungibile e se \((A,B)\) è completamente raggiungibile}

A questo punto si chiarisce bene il significato pratico delle due forme.

\begin{itemize}
    \item Se \((A,B)\) \textbf{non è completamente raggiungibile}, la forma utile è la
    \[
    \boxed{\text{forma standard di raggiungibilità}}
    \]
    perché separa esplicitamente la parte raggiungibile dalla parte non raggiungibile.

    \item Se \((A,B)\) \textbf{è completamente raggiungibile}, la forma più utile è la
    \[
    \boxed{\text{forma canonica di controllo}}
    \]
    perché permette di leggere direttamente il polinomio caratteristico e di progettare la retroazione in modo naturale.
\end{itemize}

%=========================================================
\section{Forma canonica di controllo}

\subsection{Punto di partenza: un’equazione scalare di ordine \(n\)}

Consideriamo una sequenza scalare \(z(k)\) che soddisfa
\begin{equation}
D^n z(k)+\sum_{i=0}^{n-1} a_i D^i z(k)=u(k),
\label{eq:L14_eq_scalare}
\end{equation}
dove \(D\) è l’operatore di avanzamento:
\[
Dz(k)=z(k+1),\qquad D^2z(k)=z(k+2),\quad \dots
\]

Introduciamo lo stato
\[
x_1=z,\qquad
x_2=Dz,\qquad
x_3=D^2z,\qquad \dots,\qquad
x_n=D^{n-1}z.
\]
Allora
\[
x_1^+=x_2,\qquad
x_2^+=x_3,\qquad
\dots,\qquad
x_{n-1}^+=x_n.
\]
Inoltre, dalla \eqref{eq:L14_eq_scalare},
\[
x_n^+=D^n z
=
-\sum_{i=0}^{n-1} a_i D^i z + u
=
-a_0x_1-a_1x_2-\cdots-a_{n-1}x_n+u.
\]

Quindi il sistema in forma di stato è
\[
x^+=A_cx+B_cu,
\]
con
\[
A_c=
\begin{bmatrix}
0 & 1 & 0 & \cdots & 0\\
0 & 0 & 1 & \cdots & 0\\
\vdots & \vdots & \vdots & \ddots & \vdots\\
0 & 0 & 0 & \cdots & 1\\
-a_0 & -a_1 & -a_2 & \cdots & -a_{n-1}
\end{bmatrix},
\qquad
B_c=
\begin{bmatrix}
0\\0\\\vdots\\0\\1
\end{bmatrix}.
\]

Questa è la \textbf{forma canonica di controllo}.

\subsection{Uscita associata}

Se l’uscita è della forma
\[
y(k)=\sum_{j=0}^{p} b_j D^j z(k),
\qquad p\le n-1,
\]
allora, per come sono state definite le variabili di stato,
\[
y(k)=b_0x_1(k)+b_1x_2(k)+\cdots+b_p x_{p+1}(k).
\]
Quindi
\[
C_c=
\begin{bmatrix}
b_0 & b_1 & \dots & b_p & 0 & \dots & 0
\end{bmatrix},
\qquad
D_c=0.
\]

\subsection{Polinomio caratteristico}

La scelta dei coefficienti nell’ultima riga di \(A_c\) non è casuale.  
Infatti
\[
p(\lambda)=\det(\lambda I-A_c)
=
\lambda^n+a_{n-1}\lambda^{n-1}+\cdots+a_1\lambda+a_0.
\]
Quindi i coefficienti del polinomio caratteristico compaiono direttamente in \(A_c\).

%=========================================================
\section{La forma canonica di controllo è sempre completamente raggiungibile}

Calcoliamo la matrice di raggiungibilità della forma canonica di controllo:
\[
R_c=
\begin{bmatrix}
B_c & A_cB_c & A_c^2B_c & \cdots & A_c^{n-1}B_c
\end{bmatrix}.
\]

Questa matrice ha sempre rango massimo.  
Infatti, osservando le colonne:
\begin{itemize}
    \item la prima è \(e_n\);
    \item la seconda contiene \(e_{n-1}\) più combinazioni delle precedenti;
    \item la terza contiene \(e_{n-2}\) più combinazioni delle precedenti;
    \item e così via.
\end{itemize}

Con una semplice eliminazione si vede che le colonne sono linearmente indipendenti. Quindi
\[
\det(R_c)\neq 0,
\qquad
\operatorname{rank}(R_c)=n.
\]

Conclusione:
\[
\boxed{\text{la forma canonica di controllo è sempre completamente raggiungibile.}}
\]

%=========================================================
\section{Come costruire il cambio di base verso la forma canonica di controllo}

Se il sistema \((A,B)\) è completamente raggiungibile, allora può essere portato in forma canonica di controllo.

Sia
\[
R=\begin{bmatrix}B & AB & \cdots & A^{n-1}B\end{bmatrix}
\]
la matrice di raggiungibilità del sistema dato, e sia
\[
R_c
\]
la matrice di raggiungibilità della forma canonica di controllo avente lo stesso polinomio caratteristico.

Poiché entrambe sono invertibili, il cambio di base si può scegliere come
\begin{equation}
T=RR_c^{-1}.
\label{eq:L14_T_RRc}
\end{equation}

Allora, ponendo \(x=Tz\), si ottiene
\[
A_c=T^{-1}AT,\qquad
B_c=T^{-1}B.
\]

\subsection{Perché questa formula funziona}

L’idea è semplice:
\begin{itemize}
    \item \(R\) descrive la catena dei vettori di raggiungibilità del sistema originale;
    \item \(R_c\) descrive la stessa catena per la forma canonica.
\end{itemize}

Moltiplicare per \(R_c^{-1}\) e poi per \(R\) significa costruire il cambio di coordinate che manda la base canonica della forma di controllo nella base di raggiungibilità del sistema originario.

\subsection{Esempio}

Consideriamo
\[
A=
\begin{bmatrix}
-1 & 1 & 0\\
0 & -1 & 1\\
0 & 0 & 2
\end{bmatrix},
\qquad
B=
\begin{bmatrix}
0\\1\\1
\end{bmatrix}.
\]
Il polinomio caratteristico è
\[
p(\lambda)=(\lambda+1)^2(\lambda-2)=\lambda^3-3\lambda-2.
\]
Quindi la corrispondente forma canonica di controllo è
\[
A_c=
\begin{bmatrix}
0 & 1 & 0\\
0 & 0 & 1\\
2 & 3 & 0
\end{bmatrix},
\qquad
B_c=
\begin{bmatrix}
0\\0\\1
\end{bmatrix}.
\]

La matrice di raggiungibilità del sistema originale è
\[
R=
\begin{bmatrix}
0 & 1 & -1\\
1 & 0 & 2\\
1 & 2 & 4
\end{bmatrix},
\]
mentre
\[
R_c=
\begin{bmatrix}
0 & 0 & 1\\
0 & 1 & 0\\
1 & 0 & 3
\end{bmatrix}.
\]
Da cui
\[
R_c^{-1}=
\begin{bmatrix}
-3 & 0 & 1\\
0 & 1 & 0\\
1 & 0 & 0
\end{bmatrix}.
\]
Quindi
\[
T=RR_c^{-1}
=
\begin{bmatrix}
-1 & 1 & 0\\
-1 & 0 & 1\\
-1 & 2 & 1
\end{bmatrix}.
\]

Questo è il cambio di base che porta il sistema in forma canonica di controllo.

%=========================================================
\section{Retroazione di stato e autovalori non raggiungibili}

Supponiamo ora di avere il sistema già in forma standard di raggiungibilità:
\[
\begin{cases}
z_R^+=A_R z_R + A_{RN} z_N + B_R u,\\
z_N^+=A_N z_N.
\end{cases}
\]

Se imponiamo una retroazione di stato
\[
u=Kx,
\qquad
K=\begin{bmatrix}K_R & K_N\end{bmatrix},
\]
allora la matrice di stato a ciclo chiuso diventa
\[
A+BK
=
\begin{bmatrix}
A_R & A_{RN}\\
0 & A_N
\end{bmatrix}
+
\begin{bmatrix}
B_R\\0
\end{bmatrix}
\begin{bmatrix}
K_R & K_N
\end{bmatrix}
=
\begin{bmatrix}
A_R+B_RK_R & A_{RN}+B_RK_N\\
0 & A_N
\end{bmatrix}.
\]

La cosa fondamentale è che il blocco \(A_N\) \textbf{non cambia}.  
Quindi gli autovalori della parte non raggiungibile restano invariati sotto retroazione di stato.

\begin{quote}
In altre parole: con una retroazione di stato puoi spostare solo gli autovalori raggiungibili.
\end{quote}

Questo è uno dei motivi più profondi per cui ha senso distinguere la parte raggiungibile da quella non raggiungibile.

%=========================================================
\section{Introduzione all’osservabilità}

Passiamo ora al problema duale.

\subsection{Domanda fondamentale}

Dato il modello del sistema
\[
(A,B,C,D),
\]
e supponendo noti:
\begin{itemize}
    \item l’ingresso \(u(\cdot)\),
    \item l’uscita misurata \(y(\cdot)\),
\end{itemize}
si possono ricostruire informazioni sullo stato iniziale \(x(0)\)?  
E più in generale sullo stato \(x(t)\)?

Questo è il \textbf{problema di osservazione dello stato}.

\subsection{Stati indistinguibili}

Due stati iniziali \(x_0\) e \(\bar x_0\) si dicono \textbf{indistinguibili} in un intervallo \([0,t_f]\) se, a parità di ingresso, producono la stessa uscita per tutto l’intervallo:
\[
y(\tau; x_0,u)=y(\tau; \bar x_0,u)
\qquad
\forall \tau\in[0,t_f],\ \forall u.
\]

Indichiamo con
\[
\mathcal N_{t_f}(x_0)
=
\left\{
\bar x_0\in\mathbb R^n:
y(\tau; x_0,u)=y(\tau; \bar x_0,u)\ \ \forall \tau,\forall u
\right\}
\]
l’insieme degli stati indistinguibili da \(x_0\).

Il sistema è \textbf{completamente osservabile} se l’unico stato indistinguibile da \(x_0\) è \(x_0\) stesso, cioè se l’unico stato indistinguibile dall’origine è l’origine:
\[
\mathcal N_{t_f}(0)=\{0\}.
\]

%=========================================================
\section{Caso continuo: formula esplicita del sottospazio di inosservabilità}

Per il sistema continuo
\[
\dot x=Ax+Bu,\qquad y=Cx+Du,
\]
con stato iniziale \(x(0)=x_0\), si ha
\[
y(\tau)=Ce^{A\tau}x_0+\int_0^\tau Ce^{A(\tau-\sigma)}Bu(\sigma)\,d\sigma + Du(\tau).
\]

Se \(x_0\) e \(\bar x_0\) sono indistinguibili, allora, sottraendo le due uscite a parità di ingresso, i termini forzati si cancellano e rimane
\[
Ce^{A\tau}(x_0-\bar x_0)=0
\qquad
\forall \tau\in[0,t_f].
\]

Dunque
\[
x_0-\bar x_0 \in \ker\!\big(Ce^{A\tau}\big)\qquad \forall \tau\in[0,t_f].
\]
Ne segue che
\begin{equation}
\mathcal N_{t_f}(0)=\bigcap_{\tau\in[0,t_f]} \ker\!\big(Ce^{A\tau}\big).
\label{eq:L14_Ntf}
\end{equation}

Questo è il \textbf{sottospazio di inosservabilità}.

%=========================================================
\section{Operatore di osservazione e gramiano di osservabilità}

Definiamo l’operatore
\[
\psi : \mathbb R^n \to L^2([0,t_f],\mathbb R^p),
\qquad
\psi(x)=Ce^{A(\cdot)}x.
\]
Questo operatore associa a uno stato iniziale la corrispondente \emph{uscita libera}.

L’operatore aggiunto soddisfa
\[
\langle \psi(x), y\rangle=\langle x,\psi^\ast(y)\rangle,
\]
e si calcola come
\[
\psi^\ast(y)=\int_0^{t_f} e^{A^\top \tau} C^\top y(\tau)\,d\tau.
\]

La composizione \(\psi^\ast\psi\) è allora
\[
\psi^\ast\psi(x)
=
\int_0^{t_f} e^{A^\top \tau}C^\top C e^{A\tau}\,d\tau \, x.
\]
Introduciamo quindi il \textbf{gramiano di osservabilità}
\begin{equation}
G_O(0,t_f)=\int_0^{t_f} e^{A^\top \tau}C^\top C e^{A\tau}\,d\tau.
\label{eq:L14_gramiano_oss}
\end{equation}

Poiché
\[
\psi^\ast\psi=G_O(0,t_f),
\]
vale
\[
\ker(\psi)=\ker(G_O(0,t_f)).
\]
Ma \(\ker(\psi)\) coincide proprio con il sottospazio degli stati che non lasciano traccia in uscita. Quindi
\[
\boxed{
\mathcal N_{t_f}(0)=\ker(G_O(0,t_f)).
}
\]

Ne segue che il sistema è completamente osservabile se e solo se
\[
G_O(0,t_f)\ \text{è invertibile}.
\]

%=========================================================
\section{Matrice di osservabilità e parallelismo con la raggiungibilità}

Come per la raggiungibilità, esiste anche una caratterizzazione puramente algebrica.

Definiamo la \textbf{matrice di osservabilità}
\[
\mathcal O=
\begin{bmatrix}
C\\
CA\\
CA^2\\
\vdots\\
CA^{n-1}
\end{bmatrix}.
\]

Allora vale il parallelo perfetto con la parte precedente:
\[
\ker(G_O(0,t_f))=\ker(\mathcal O).
\]
Quindi il sistema è completamente osservabile se e solo se
\[
\operatorname{rank}(\mathcal O)=n.
\]

\bigskip

Confronto diretto:

\[
\text{raggiungibilità}
\quad\longleftrightarrow\quad
R=\begin{bmatrix}B & AB & \cdots & A^{n-1}B\end{bmatrix},
\]
\[
\text{osservabilità}
\quad\longleftrightarrow\quad
\mathcal O=
\begin{bmatrix}
C\\CA\\ \vdots\\ CA^{n-1}
\end{bmatrix}.
\]

Questo forte parallelismo è il preludio alla \textbf{dualità} tra raggiungibilità e osservabilità.

%=========================================================
\section{Conclusione della lezione}

In questa lezione abbiamo visto quattro idee fondamentali.

\begin{enumerate}
    \item La \textbf{forma standard di raggiungibilità} separa lo stato in una parte raggiungibile e una non raggiungibile.

    \item La \textbf{forma canonica di controllo} è il modello di riferimento dei sistemi completamente raggiungibili:
    \[
    A_c=
    \begin{bmatrix}
    0 & 1 & 0 & \cdots & 0\\
    0 & 0 & 1 & \cdots & 0\\
    \vdots & \vdots & \vdots & \ddots & \vdots\\
    0 & 0 & 0 & \cdots & 1\\
    -a_0 & -a_1 & -a_2 & \cdots & -a_{n-1}
    \end{bmatrix},
    \qquad
    B_c=e_n.
    \]

    \item Gli \textbf{autovalori non raggiungibili} non possono essere spostati tramite retroazione di stato.

    \item L’\textbf{osservabilità} è il problema duale: capire se dallo storico di ingressi e uscite è possibile ricostruire lo stato.
\end{enumerate}

\bigskip

La prossima lezione potrà riprendere naturalmente da qui, sviluppando:
\begin{itemize}
    \item la forma standard di osservabilità;
    \item la dualità con la raggiungibilità;
    \item gli osservatori di stato.
\end{itemize}